% =========================================================== 1. cover
{\setbeamertemplate{footline}{}
\begin{frame}[plain]
  \vspace{1.7em}
  {\color{brick}\rule{3.6em}{2.4pt}}

  \vspace{1.1em}
  {\fontsize{29}{33}\selectfont\bfseries When the price tells the truth\par}

  \vspace{0.55em}
  {\fontsize{14.5}{18}\selectfont\color{muted} One battery, two hours, and the value of flexibility\par}

  \vspace{1.5em}
  \begin{minipage}{0.70\textwidth}
    \fontsize{9.6}{13}\selectfont
    A minimal example, with two generators and one ramp constraint, showing where
    a battery's revenue comes from and why it is exactly equal to what the battery
    saves the system. Then the same reading on a real day of the Brazilian system.
  \end{minipage}

  \vfill
  {\fontsize{8.4}{10}\selectfont\color{muted} Alexandre Street\ \ \textbullet\ \ August 2026}
  \vspace{0.8em}
\end{frame}}

% =========================================================== 2. the problem
\begin{frame}
  \framesubtitle{The problem}
  \frametitle{Demand rises 30 MWh in one hour. The cheap unit rises 25.}

  \begin{columns}[T,onlytextwidth]
  \begin{column}{0.42\textwidth}
    \vspace{0.9em}
    \fontsize{9.4}{12.6}\selectfont
    \begin{itemize}\setlength{\itemsep}{0.5em}
      \item Two thermal units, unlimited capacity.
      \item Variable cost of \textbf{100} and \textbf{300} R\$/MWh.
      \item Ramp-up limit of \textbf{25 MWh/h}, the same for both.
      \item Demand of \textbf{100 MWh} at t1 and \textbf{130 MWh} at t2.
    \end{itemize}

    \vspace{0.9em}
    \gloss{t1 is the first hour of the horizon, so the ramp constrains only the
    step from t1 to t2.}
  \end{column}

  \begin{column}{0.55\textwidth}
    \vspace{0.7em}
    \centering
    \begin{tikzpicture}[y=0.088cm,x=0.86cm]
      \draw[rulec,line width=0.7pt] (-0.30,90) -- (-0.30,136);
      \foreach \v in {100,125,130}{
        \draw[rulec,line width=0.7pt] (-0.42,\v) -- (-0.30,\v);
        \node[left,muted,font=\fontsize{6.4}{7}\selectfont] at (-0.46,\v) {\v};
      }
      \draw[ink,line width=2.0pt] (0.10,100) -- (1.55,100);
      \node[above,font=\fontsize{7.6}{8}\selectfont\bfseries,ink] at (0.82,101.5) {100 MWh};
      \draw[ink,line width=2.0pt] (2.75,130) -- (4.20,130);
      \node[above,font=\fontsize{7.6}{8}\selectfont\bfseries,ink] at (3.47,131.5) {130 MWh};
      \draw[brick,line width=1.0pt,dashed] (2.75,125) -- (4.20,125);
      \node[below,font=\fontsize{6.4}{7}\selectfont,brick] at (3.47,123.6) {125 = 100 + 25};
      \draw[brick,line width=0.9pt] (4.42,125) -- (4.42,130);
      \draw[brick,line width=0.9pt] (4.34,125) -- (4.50,125);
      \draw[brick,line width=0.9pt] (4.34,130) -- (4.50,130);
      \node[right,font=\fontsize{7.4}{8.4}\selectfont\bfseries,brick,align=left]
        at (4.50,127.5) {5 MWh\\[-1pt]{\fontsize{6.2}{7.4}\selectfont\mdseries expensive unit}};
      \draw[->,>=stealth,muted,line width=0.9pt] (1.62,100.8) -- (2.68,124.2);
      \node[muted,font=\fontsize{6.4}{7.6}\selectfont,anchor=south east,align=right]
        at (2.30,110) {ramp limit\\[-1pt]25 MWh/h};
      \node[muted,font=\fontsize{7.4}{8}\selectfont] at (0.82,93.5) {t1};
      \node[muted,font=\fontsize{7.4}{8}\selectfont] at (3.47,93.5) {t2};
      \draw[rulec,line width=0.7pt] (-0.30,90) -- (5.35,90);
    \end{tikzpicture}

    \vspace{0.5em}
    \gloss{\centering Vertical axis truncated at 90 MWh.}
  \end{column}
  \end{columns}
\end{frame}

% =========================================================== 3. case A
\begin{frame}
  \framesubtitle{Case A \ \textbullet\ \ no battery}
  \frametitle{The optimal dispatch costs R\$ 24,000, and the price at t1 turns negative.}

  \begin{columns}[T,onlytextwidth]
  \begin{column}{0.635\textwidth}
    \vspace{0.9em}
    {\tabfont
    \setlength{\tabcolsep}{4.0pt}
    \renewcommand{\arraystretch}{1.19}
    \begin{tabular}{l r r @{\hspace{1.1em}} r r r}
      \toprule
       & \textbf{t1} & \textbf{t2} & \textbf{Revenue} & \textbf{Cost} & \textbf{Profit} \\
      \midrule
      Demand         & 100 & 130 & $-$29,000 & & \\
      Thermal 1      & 100 & 125 & 27,500 & 22,500 & 5,000 \\
      Thermal 2      & 0   & 5   & 1,500  & 1,500  & 0 \\
      \addlinespace[0.1em]\midrule
      Hourly cost    & 10,000 & 14,000 & & & \\
      Spot price     & \hl{$-$100} & \textbf{300} & & & \\
      \addlinespace[0.1em]\midrule
      \textbf{Total} & & & \textbf{0} & \textbf{24,000} & \textbf{5,000} \\
      \bottomrule
    \end{tabular}}

    \vspace{0.8em}
    \gloss{Quantities in MWh; revenue, cost, profit and hourly cost in R\$; spot
    price in R\$/MWh. Demand's revenue enters negative, because it is what
    consumers pay. Each hour's spot price is the marginal operating cost: the
    change in total cost over the two periods when 1 MWh of demand is added in
    that hour and the other one is held fixed.}
  \end{column}

  \begin{column}{0.335\textwidth}
    \vspace{0.9em}
    \panel{%
      \fontsize{8.2}{10.6}\selectfont
      \textbf{The dispatch model}\par
      \vspace{0.55em}
      $\displaystyle\min_{g\,\ge\,0}\ \sum_{t}\sum_{i} c_i\,g_{i,t}$\par
      \vspace{0.55em}
      subject to\par
      \vspace{0.30em}
      $\textstyle\sum_i g_{i,t} = d_t \quad (\lambda_t)$\par
      \vspace{0.42em}
      $g_{i,2}-g_{i,1} \le R \quad (\mu_i)$\par
      \vspace{0.62em}
      \gloss{$\lambda_t$ is the price of energy in hour $t$. $\mu_i$ is what one
      more MWh of ramp is worth to the system.}
    }
  \end{column}
  \end{columns}
\end{frame}

% =========================================================== 4. duals
\begin{frame}
  \framesubtitle{Reading the duals}
  \frametitle{The negative price is not a flaw in the model. It is the right arithmetic.}

  \begin{columns}[T,onlytextwidth]
  \begin{column}{0.50\textwidth}
    \vspace{0.5em}
    {\fontsize{9.6}{12}\selectfont\hspace{0.9em}
    $\begin{aligned}
      \mu_1 &= c_2-c_1        &&= \hlm{200} \\[3pt]
      \lambda_1 &= c_1-\mu_1  &&= \hlm{-100} \\[3pt]
      \lambda_2 &= c_1+\mu_1  &&= \hlm{+300} \\[3pt]
      \lambda_2-\lambda_1 &= 2\,(c_2-c_1) &&= \hgm{400}
    \end{aligned}$\par}

    \vspace{0.85em}
    {\fontsize{8.6}{11.4}\selectfont
    \begin{itemize}\setlength{\itemsep}{0.34em}
      \item Ramp headroom is worth the difference between the two units.
      \item Generating at t1 costs 100 and \emph{creates} headroom worth 200. Net, $-$100.
      \item Generating at t2 costs 100 and \emph{consumes} headroom worth 200. Total, 300.
      \item The spread between the hours is twice the cost difference.
    \end{itemize}}
  \end{column}

  \begin{column}{0.47\textwidth}
    \vspace{0.5em}
    \panel{%
      \fontsize{8.6}{11}\selectfont
      \textbf{Revenue adequacy}\par
      \vspace{0.15em}
      Consumers pay 29,000 and generators receive 29,000. The revenue column sums
      to zero.

      \vspace{0.8em}
      \textbf{Cost recovery}\par
      \vspace{0.15em}
      Thermal 1 sells at $-$100 at t1 and still closes with a profit of 5,000,
      which is exactly $\mu_1 \times R = 200 \times 25$, the rent on its own ramp.
      Thermal 2 closes at zero. No generator loses money.
    }

    \vspace{0.7em}
    \gloss{Electricity has no waste bin. Since it is impossible to generate more
    than is consumed, the negative price at t1 is the system saying that consuming
    in that hour is worth 100 R\$/MWh.}
  \end{column}
  \end{columns}
\end{frame}

% =========================================================== 5. case B
\begin{frame}
  \framesubtitle{Case B \ \textbullet\ \ with a 1 \nocaps{MWh} battery}
  \frametitle{The battery buys at $-$100, sells at $+$300, and earns 400.}

  \begin{columns}[T,onlytextwidth]
  \begin{column}{0.635\textwidth}
    \vspace{0.9em}
    {\tabfont
    \setlength{\tabcolsep}{4.0pt}
    \renewcommand{\arraystretch}{1.19}
    \begin{tabular}{l r r @{\hspace{1.1em}} r r r}
      \toprule
       & \textbf{t1} & \textbf{t2} & \textbf{Revenue} & \textbf{Cost} & \textbf{Profit} \\
      \midrule
      Demand         & 100 & 130 & $-$29,000 & & \\
      Thermal 1      & 101 & 126 & 27,700 & 22,700 & 5,000 \\
      Thermal 2      & 0   & 3   & 900    & 900    & 0 \\
      Battery        & \hg{$-$1} & \hg{$+$1} & \hg{400} & \hg{0} & \hg{400} \\
      \addlinespace[0.1em]\midrule
      Hourly cost    & 10,100 & 13,500 & & & \\
      Spot price     & \hl{$-$100} & \textbf{300} & & & \\
      \addlinespace[0.1em]\midrule
      \textbf{Total} & & & \textbf{0} & \textbf{23,600} & \textbf{5,400} \\
      \bottomrule
    \end{tabular}}

    \vspace{0.8em}
    \gloss{The battery charges 1 MWh at t1, when the price is $-$100, and is
    therefore paid 100 to consume. It discharges that same 1 MWh at t2 and sells
    at 300. Revenue of 400 and zero cost. Prices do not change, because the units
    at the margin are still the same.}
  \end{column}

  \begin{column}{0.335\textwidth}
    \vspace{0.9em}
    \panel{%
      \fontsize{8.2}{10.6}\selectfont
      \textbf{The model with the battery}\par
      \vspace{0.55em}
      $\displaystyle\min_{g,\,e\,\ge\,0}\ \sum_{t}\sum_{i} c_i\,g_{i,t}$\par
      \vspace{0.55em}
      subject to\par
      \vspace{0.30em}
      $\textstyle\sum_i g_{i,1} = d_1 + e \quad (\lambda_1)$\par
      \vspace{0.42em}
      $\textstyle\sum_i g_{i,2} + e = d_2 \quad (\lambda_2)$\par
      \vspace{0.42em}
      $g_{i,2}-g_{i,1} \le R \quad (\mu_i)$\par
      \vspace{0.42em}
      $e \le E \quad (\beta)$\par
      \vspace{0.62em}
      \gloss{$e$ is the energy shifted from t1 to t2, with $E = 1$ MWh and 100\%
      efficiency. The dual of the capacity limit is the price difference,
      \mbox{$\beta = 400$}.}
    }
  \end{column}
  \end{columns}
\end{frame}

% =========================================================== 6. comparison
\begin{frame}
  \framesubtitle{The two cases side by side}
  \frametitle{The battery's profit is the saving it creates for the system.}

  \vspace{0.7em}
  \centering
  {\tabfont
  \setlength{\tabcolsep}{5pt}
  \renewcommand{\arraystretch}{1.26}
  \begin{tabular}{l r r r @{\hspace{1.5em}} r r r}
    \toprule
     & \multicolumn{3}{c}{\textbf{Without battery}} & \multicolumn{3}{c}{\textbf{With battery}} \\
    \cmidrule(r){2-4}\cmidrule(l){5-7}
     & Revenue & Cost & Profit & Revenue & Cost & Profit \\
    \midrule
    Thermal 1     & 27,500 & 22,500 & 5,000 & 27,700 & 22,700 & 5,000 \\
    Thermal 2     & 1,500  & 1,500  & 0     & 900    & 900    & 0 \\
    Battery       &        &        &       & \hg{400} & \hg{0} & \hg{400} \\
    Consumers     & $-$29,000 & & & $-$29,000 & & \\
    \addlinespace[0.1em]\midrule
    \textbf{Total} & \textbf{0} & \textbf{24,000} & \textbf{5,000}
                   & \textbf{0} & \textbf{23,600} & \textbf{5,400} \\
    \bottomrule
  \end{tabular}}

  \vspace{1.2em}
  \begin{columns}[T,onlytextwidth]
  \begin{column}{0.305\textwidth}
    \panel{\centering\fontsize{8.2}{10.4}\selectfont
      Operating saving\par\vspace{0.25em}
      {\fontsize{11}{12}\selectfont\bfseries 24,000 $\rightarrow$ 23,600}\par
      \vspace{0.15em}\hg{$-$400 R\$}}
  \end{column}
  \begin{column}{0.305\textwidth}
    \panel{\centering\fontsize{8.2}{10.4}\selectfont
      Battery profit\par\vspace{0.25em}
      {\fontsize{11}{12}\selectfont\bfseries 400 $-$ 0}\par
      \vspace{0.15em}\hg{$+$400 R\$}}
  \end{column}
  \begin{column}{0.305\textwidth}
    \panel{\centering\fontsize{8.2}{10.4}\selectfont
      Paid by consumers\par\vspace{0.25em}
      {\fontsize{11}{12}\selectfont\bfseries 29,000 $=$ 29,000}\par
      \vspace{0.15em}\hg{unchanged}}
  \end{column}
  \end{columns}
\end{frame}

% =========================================================== 7. saturation
\begin{frame}
  \framesubtitle{How much to build}
  \frametitle{The battery's income is a scarcity rent, and it vanishes at 2.5 MWh.}

  \begin{columns}[T,onlytextwidth]
  \begin{column}{0.60\textwidth}
    \vspace{1.1em}
    {\tabfont
    \setlength{\tabcolsep}{4.0pt}
    \renewcommand{\arraystretch}{1.26}
    \begin{tabular}{l r r r r}
      \toprule
      \textbf{Installed} & \textbf{Total} & \textbf{Price} & \textbf{Price}
        & \textbf{Marginal} \\
      \textbf{battery} & \textbf{cost} & \textbf{t1} & \textbf{t2}
        & \textbf{value} \\
      \midrule
      0 MWh             & 24,000 & $-$100 & 300 & 400 \\
      1 MWh             & 23,600 & $-$100 & 300 & 400 \\
      2 MWh             & 23,200 & $-$100 & 300 & 400 \\
      \addlinespace[0.12em]
      2.5 MWh or more   & 23,000 & \hg{100} & \hg{100} & \hg{0} \\
      \bottomrule
    \end{tabular}}

    \vspace{0.8em}
    \gloss{Total cost in R\$; prices and marginal value in R\$/MWh. The marginal
    value is what the next MWh of battery is worth to the system.}
  \end{column}

  \begin{column}{0.37\textwidth}
    \vspace{1.1em}
    {\fontsize{9.0}{12.0}\selectfont
    With 2.5 MWh of storage the ramp stops binding, the expensive unit leaves the
    dispatch, and both hours settle at 100 R\$/MWh. From there on, the next MWh of
    battery is worth zero.\par}

    \vspace{1.25em}
    \thesis{\fontsize{8.8}{11.4}\selectfont
      The same price that pays the first battery tells us when the next one stops
      being needed. That is how price connects the market, and the incentives it
      creates, to the needs of those who operate the system.}
  \end{column}
  \end{columns}
\end{frame}
